\documentclass[11pt,a4paper]{article}
\usepackage[utf8]{inputenc}
\usepackage[T1]{fontenc}
\usepackage{lmodern}
\usepackage[portuguese]{babel}
\usepackage[a4paper,margin=2.5cm]{geometry}
\usepackage{xcolor}
\definecolor{TextEspresso}{HTML}{4A4238}
\definecolor{NudeTaupe}{HTML}{BFAFA6}
\definecolor{SoftBlush}{HTML}{D9C5B2}
\usepackage{titlesec}
\usepackage{enumitem}
\usepackage{listings}
\usepackage{hyperref}
\hypersetup{colorlinks=true, linkcolor=TextEspresso, urlcolor=TextEspresso}

\titleformat{\section}{\color{TextEspresso}\Large\bfseries}{\thesection}{0.5em}{}[{\color{NudeTaupe}\titlerule[0.8pt]}]
\titleformat{\subsection}{\color{TextEspresso}\large\bfseries}{\thesubsection}{0.5em}{}

\lstset{
  basicstyle=\ttfamily\small,
  breaklines=true,
  frame=single,
  rulecolor=\color{NudeTaupe},
  columns=fullflexible,
  showstringspaces=false
}

\newcommand{\guiaotitle}[2]{%
  \begin{center}
  {\Large\bfseries\color{TextEspresso} #1}\\[0.3cm]
  {\normalsize\color{NudeTaupe} #2}\\[0.2cm]
  {\color{NudeTaupe}\rule{4cm}{0.5pt}}
  \end{center}
  \vspace{0.5cm}
}

\begin{document}
\guiaotitle{Guião 07 -- IaC com Ansible e Terraform no Proxmox}{Infraestruturas de Computação na Nuvem, Aula 07}

\section{Objetivos}
\begin{itemize}[itemsep=0.2cm]
  \item Configurar o Ansible e definir um \emph{inventory} com as VMs Proxmox como alvos.
  \item Escrever um \emph{playbook} idempotente que instala e configura o Nginx, com variáveis, templates e handlers.
  \item Demonstrar concretamente a idempotência, correndo o mesmo playbook duas vezes.
  \item Provisionar uma VM no Proxmox de forma declarativa com Terraform, observando o ciclo \texttt{plan}/\texttt{apply}/\texttt{destroy}.
  \item Distinguir, na prática, provisionamento (Terraform) de gestão de configuração (Ansible).
\end{itemize}

\section{Pré-requisitos e Material}
\begin{itemize}[itemsep=0.2cm]
  \item A VM \texttt{icn-lab} (Guião 01), que servirá de \textbf{máquina de controlo} -- é onde o Ansible e o Terraform correm.
  \item Pelo menos uma segunda VM no Proxmox, para servir de alvo do Ansible (pode reutilizar uma VM das aulas anteriores, ou criar uma nova).
  \item Acesso SSH da máquina de controlo para a(s) VM(s) alvo, com privilégios de \texttt{sudo}.
  \item Para a Parte 3 (Terraform): um \emph{API token} do Proxmox com permissões de criação de VMs, fornecido pela docente.
\end{itemize}

\textbf{Nota sobre a Parte 3:} a criação de VMs por Terraform exige permissões elevadas no Proxmox. Consoante a política do servidor da UC, esta parte pode ser feita por cada aluno com um token restrito a um \emph{pool} próprio, ou realizada como demonstração pela docente. Confirme antes de começar.

\section{Introdução}
Na aula teórica distinguimos duas camadas complementares de automatização: o \textbf{provisionamento} (que recursos existem) e a \textbf{gestão de configuração} (o que corre dentro deles). Neste guião vamos exercitar as duas, na ordem inversa da que seria usada em produção -- primeiro a configuração, depois o provisionamento -- por ser a progressão pedagogicamente mais simples.

Na Parte 1 e 2 usamos o \textbf{Ansible}, uma ferramenta \emph{agentless} que comunica por SSH e descreve o estado desejado em YAML. Na Parte 3 usamos o \textbf{Terraform} para criar uma VM no Proxmox a partir de código declarativo, fechando o ciclo: Terraform cria a máquina, Ansible configura-a.

Repare que, nas aulas anteriores, criámos VMs manualmente (interface web) e por comandos (\texttt{qm}). O que fazemos hoje é o passo seguinte: descrever a infraestrutura em ficheiros versionáveis, em vez de a construir por ações irrepetíveis.

\section{Parte 1 -- Ansible: preparação}
\begin{enumerate}[label=\arabic*.]

  \item \textbf{Instalar o Ansible} na máquina de controlo:
\begin{lstlisting}
$ sudo apt update
$ sudo apt install -y ansible
$ ansible --version
\end{lstlisting}

  \item \textbf{Configurar acesso SSH por chave} para a(s) VM(s) alvo, evitando pedidos de password a cada execução:
\begin{lstlisting}
$ ssh-keygen -t ed25519 -C "ansible-icn"
$ ssh-copy-id utilizador@<IP-da-VM-alvo>
$ ssh utilizador@<IP-da-VM-alvo> "hostname"
\end{lstlisting}

  \item \textbf{Criar a pasta do projeto:}
\begin{lstlisting}
$ mkdir -p ~/ansible-nginx/templates && cd ~/ansible-nginx
\end{lstlisting}

  \item \textbf{Criar o \emph{inventory}} \texttt{inventory.ini}, com as VMs alvo (substituir pelos IPs reais):
\begin{lstlisting}
[webservers]
vm1 ansible_host=<IP-da-VM-alvo-1> ansible_user=<utilizador>
vm2 ansible_host=<IP-da-VM-alvo-2> ansible_user=<utilizador>

[webservers:vars]
ansible_python_interpreter=/usr/bin/python3
\end{lstlisting}
  Se só tiver uma VM alvo disponível, inclua apenas a linha \texttt{vm1}.

  \item \textbf{Testar a ligação a todos os hosts:}
\begin{lstlisting}
$ ansible -i inventory.ini webservers -m ping
\end{lstlisting}
  Uma resposta \texttt{pong} por host confirma que o SSH e o Python estão corretamente configurados.

  \item \textbf{Observar os \emph{facts} recolhidos automaticamente} de um host:
\begin{lstlisting}
$ ansible -i inventory.ini vm1 -m setup | head -40
\end{lstlisting}
  Localize, entre outros, os seguintes campos:
\begin{lstlisting}
ansible_distribution
ansible_memtotal_mb
ansible_default_ipv4
\end{lstlisting}
  Estes valores podem ser usados em condições e em templates.

\end{enumerate}

\section{Parte 2 -- Ansible: playbook idempotente}
\begin{enumerate}[label=\arabic*.,resume]

  \item \textbf{Criar o ficheiro de variáveis} \texttt{vars.yml}:
\begin{lstlisting}
site_name: "aula-icn"
http_port: 80
welcome_message: "Servico configurado automaticamente com Ansible"
\end{lstlisting}

  \item \textbf{Criar o template Jinja2} \texttt{templates/index.html.j2}. Note o uso de \emph{facts} para incluir informação específica de cada host:
\begin{lstlisting}
<html>
  <head><title>{{ site_name }}</title></head>
  <body>
    <h1>{{ welcome_message }}</h1>
    <p>Servido por: {{ ansible_hostname }}</p>
    <p>Sistema: {{ ansible_distribution }} {{ ansible_distribution_version }}</p>
    <p>Endereco IP: {{ ansible_default_ipv4.address }}</p>
  </body>
</html>
\end{lstlisting}
  O mesmo template gera conteúdo diferente em cada máquina -- é aqui que se vê a utilidade dos facts.

  \item \textbf{Criar o playbook} \texttt{nginx.yml}:
\begin{lstlisting}
---
- name: Instalar e configurar Nginx
  hosts: webservers
  become: true
  vars_files:
    - vars.yml

  tasks:
    - name: Instalar o pacote nginx
      apt:
        name: nginx
        state: present
        update_cache: true

    - name: Implementar pagina de boas-vindas
      template:
        src: templates/index.html.j2
        dest: /var/www/html/index.html
        owner: root
        group: root
        mode: '0644'
      notify: Reiniciar nginx

    - name: Garantir que o servico esta ativo e habilitado
      service:
        name: nginx
        state: started
        enabled: true

  handlers:
    - name: Reiniciar nginx
      service:
        name: nginx
        state: restarted
\end{lstlisting}

  \item \textbf{Validar antes de aplicar.} O modo \emph{check} simula a execução sem alterar nada -- o equivalente ao \texttt{plan} do Terraform:
\begin{lstlisting}
$ ansible-playbook -i inventory.ini nginx.yml --check --diff
\end{lstlisting}

  \item \textbf{Executar o playbook} e registar o resumo final:
\begin{lstlisting}
$ ansible-playbook -i inventory.ini nginx.yml
\end{lstlisting}
  Anote os valores de \texttt{changed} e \texttt{ok} apresentados no \texttt{PLAY RECAP}.

  \item \textbf{Demonstrar a idempotência.} Execute \emph{exatamente o mesmo comando} uma segunda vez, sem alterar nada:
\begin{lstlisting}
$ ansible-playbook -i inventory.ini nginx.yml
\end{lstlisting}
  Compare o \texttt{PLAY RECAP} com o da execução anterior. Desta vez, \texttt{changed} deve ser \texttt{0}: nada foi alterado porque o estado desejado já estava atingido. Repare também que o handler \emph{não} correu -- o Nginx não foi reiniciado desnecessariamente.

  \item \textbf{Verificar o resultado:}
\begin{lstlisting}
$ curl http://<IP-da-VM-alvo-1>
\end{lstlisting}
  Se tiver duas VMs alvo, compare as duas respostas e confirme que os campos preenchidos pelos facts diferem.

  \item \textbf{Provocar drift e corrigi-lo.} Simule uma alteração manual indevida diretamente na VM alvo:
\begin{lstlisting}
$ ssh utilizador@<IP-da-VM-alvo-1>
$ sudo rm /var/www/html/index.html
$ sudo systemctl stop nginx
$ exit
\end{lstlisting}
  Volte a correr o playbook e observe que ele deteta e repõe automaticamente o estado correto:
\begin{lstlisting}
$ ansible-playbook -i inventory.ini nginx.yml
$ curl http://<IP-da-VM-alvo-1>
\end{lstlisting}
  Este é o mecanismo central da gestão de configuração: convergir sempre para o estado declarado.

  \item \textbf{Alterar a configuração e observar o handler.} Edite \texttt{welcome\_message} em \texttt{vars.yml} e volte a correr o playbook. Confirme que apenas a task \texttt{template} reporta \texttt{changed} e que, por isso, o handler é acionado e o Nginx reiniciado.

\end{enumerate}

\section{Parte 3 -- Terraform: provisionar uma VM no Proxmox}
\begin{enumerate}[label=\arabic*.,resume]

  \item \textbf{Instalar o Terraform} na máquina de controlo:
\begin{lstlisting}
$ sudo apt install -y wget gpg
$ wget -O- https://apt.releases.hashicorp.com/gpg \
    | sudo gpg --dearmor -o /usr/share/keyrings/hashicorp-archive-keyring.gpg
$ echo "deb [signed-by=/usr/share/keyrings/hashicorp-archive-keyring.gpg] https://apt.releases.hashicorp.com $(lsb_release -cs) main" \
    | sudo tee /etc/apt/sources.list.d/hashicorp.list
$ sudo apt update && sudo apt install -y terraform
$ terraform version
\end{lstlisting}

  \item \textbf{Criar a pasta do projeto Terraform:}
\begin{lstlisting}
$ mkdir ~/terraform-proxmox && cd ~/terraform-proxmox
\end{lstlisting}

  \item \textbf{Declarar o provider.} Crie \texttt{providers.tf}. O Proxmox não tem provider oficial da HashiCorp; usa-se um provider da comunidade (confirme com a docente qual está em uso e a respetiva versão):
\begin{lstlisting}
terraform {
  required_providers {
    proxmox = {
      source  = "Telmate/proxmox"
      version = "3.0.1-rc4"
    }
  }
}

provider "proxmox" {
  pm_api_url          = var.pm_api_url
  pm_api_token_id     = var.pm_api_token_id
  pm_api_token_secret = var.pm_api_token_secret
  pm_tls_insecure     = true
}
\end{lstlisting}

  \item \textbf{Declarar as variáveis.} Crie \texttt{variables.tf}, mantendo os segredos fora do código:
\begin{lstlisting}
variable "pm_api_url" {
  type = string
}

variable "pm_api_token_id" {
  type      = string
  sensitive = true
}

variable "pm_api_token_secret" {
  type      = string
  sensitive = true
}

variable "target_node" {
  type = string
}

variable "vm_template" {
  type = string
}
\end{lstlisting}
  Os valores concretos vão num ficheiro \texttt{terraform.tfvars} \textbf{que nunca deve ser commitado} num repositório:
\begin{lstlisting}
pm_api_url          = "https://<IP-do-servidor>:8006/api2/json"
pm_api_token_id     = "<utilizador>@pve!<nome-do-token>"
pm_api_token_secret = "<segredo-do-token>"
target_node         = "<nome-do-node>"
vm_template         = "<nome-do-template>"
\end{lstlisting}
  Crie também um \texttt{.gitignore} com \texttt{terraform.tfvars} e \texttt{*.tfstate*}, pelas razões discutidas na teoria.

  \item \textbf{Declarar o recurso.} Crie \texttt{main.tf}:
\begin{lstlisting}
resource "proxmox_vm_qemu" "vm_iac" {
  name        = "icn-terraform"
  target_node = var.target_node
  clone       = var.vm_template
  full_clone  = true

  cores  = 2
  memory = 2048

  network {
    id     = 0
    model  = "virtio"
    bridge = "vmbr0"
  }
}
\end{lstlisting}
  \emph{Nota:} a sintaxe exata dos blocos varia entre versões do provider. Se o \texttt{terraform validate} acusar erros, consulte a documentação da versão instalada -- este ajuste faz parte do exercício.

  \item \textbf{Inicializar e validar:}
\begin{lstlisting}
$ terraform init
$ terraform validate
$ terraform fmt
\end{lstlisting}

  \item \textbf{Executar o \texttt{plan}.} Este é o passo conceptualmente mais importante: mostra o que \emph{seria} feito, sem fazer nada:
\begin{lstlisting}
$ terraform plan
\end{lstlisting}
  Leia com atenção o resumo final (\texttt{Plan: 1 to add, 0 to change, 0 to destroy}) e registe-o.

  \item \textbf{Aplicar:}
\begin{lstlisting}
$ terraform apply
\end{lstlisting}
  Confirme na interface web do Proxmox que a VM foi efetivamente criada.

  \item \textbf{Inspecionar o estado registado:}
\begin{lstlisting}
$ terraform state list
$ terraform show
\end{lstlisting}

  \item \textbf{Observar a deteção de drift.} Na interface web do Proxmox, altere manualmente a memória da VM criada (ex.\ para 1024~MB). De volta ao terminal:
\begin{lstlisting}
$ terraform plan
\end{lstlisting}
  O Terraform deteta a divergência entre o estado real e o declarado, e propõe repô-la. Este é o comportamento discutido na teoria a propósito do \emph{state file}.

  \item \textbf{Destruir a infraestrutura criada:}
\begin{lstlisting}
$ terraform destroy
\end{lstlisting}
  Confirme no Proxmox que a VM desapareceu.

\end{enumerate}

\section{Entregável / Questões de Verificação}
\begin{enumerate}[label=\alph*)]
  \setlength\itemsep{0.2cm}
  \item Submeta \texttt{nginx.yml}, \texttt{vars.yml}, o template Jinja2 e o \texttt{inventory.ini} (com IPs anonimizados, se necessário).
  \item Submeta os dois \texttt{PLAY RECAP} obtidos (primeira e segunda execução do playbook). Explique a diferença nos valores de \texttt{changed} recorrendo ao conceito de idempotência.
  \item Porque é que o handler não foi executado na segunda passagem, mas foi executado quando alterou o \texttt{welcome\_message}?
  \item Descreva o que aconteceu quando apagou manualmente o \texttt{index.html} e voltou a correr o playbook. Que problema operacional (visto na teoria) é assim resolvido?
  \item Submeta o resultado de \texttt{terraform plan} antes do \texttt{apply}. Que vantagem tem esta separação entre planear e aplicar, comparada com executar um script diretamente?
  \item Alterou a memória da VM manualmente no Proxmox e o Terraform detetou-o. Que ficheiro lhe permitiu detetar essa divergência, e porque é que esse ficheiro deve ser tratado como um segredo?
  \item Neste guião usou duas ferramentas distintas. Qual delas fez \emph{provisionamento} e qual fez \emph{gestão de configuração}? Descreva, em duas ou três frases, como as encadearia num fluxo de trabalho real.
  \item O ficheiro \texttt{terraform.tfvars} não deve ser commitado. Que alternativas existem para gerir segredos num projeto de IaC versionado em Git?
\end{enumerate}

\end{document}
